\section{Datasets and Experimental Methodology}
\label{sec:methodology}

The experiments test whether spatial supervision helps models learn from
small labelled datasets under limited training budgets. They cover two
medical datasets and a geological application, using both models trained
from scratch and pretrained models adapted to the task. Within each model,
training with constraints is compared with a baseline using the same
training configuration and examples. Comparisons between models also
reflect their different configurations. This chapter describes the data,
the training procedures and the measures used to assess the predictions.

\subsection{Design and Units of Comparison}
\label{sec:experimental-structure}

A run trains and evaluates one model with a specified objective. Its
labelled training examples form the \emph{support set}. A fold identifies
a partition into training and validation examples, with validation held
out from parameter fitting. A seed initializes the pseudorandom generator
controlling training randomness. Repeating seeds within a fold tests
sensitivity to initialization and data presentation; using several folds
also tests different data partitions. The dataset descriptions below
specify where support examples overlap between folds.

Training from scratch is limited by floating-point operations (FLOPs),
which measure arithmetic work; one GFLOP is one billion operations.
Medical budgets cover optimization, with calibration---the preliminary
measurement used to set constraint weights---recorded separately. The
geological budget includes both. Adaptation instead fixes the number
of updates to the pretrained model's adapters. These budgets support
comparisons within each study and do not make all six studies equal in cost.

Table~\ref{tab:experimental-design} records the six studies, including two
geological and four medical comparisons. A condition is a particular model,
dataset and objective. Each condition has nine fold/seed cells, one run for
each combination of three folds and three seeds. Evaluation uses the
final saved model state, or \emph{checkpoint},
allowed by each training budget. Runs do not stop early or replace that
state with a checkpoint chosen for its validation score.

Throughout the comparisons, \emph{Baseline} denotes the supervised training
objective configured for the particular model and dataset. \emph{Baseline}
$+B$, for example, retains that objective and adds boundary-band
supervision; the same convention applies to combinations of spatial terms.
The baseline may itself combine several losses, as specified in the
model-specific protocols below. In tables and charts, $B$, $E$, $S$ and $D$
abbreviate boundary bands, edges, signed distance and depth supervision,
respectively.

\begin{table}[htbp]
\centering\small\setstretch{1.05}
\begin{tabular}{@{}P{2.5cm}P{3.0cm}P{3.2cm}P{5.7cm}@{}}
\hline
Study & Support per fold & Folds; seeds & Conditions and budget \\
\hline
GeoFWI3D, from scratch & 19 scenes & 0,1,2; 0,1,2 & 4 models $\times$ 14 objectives; 504 runs; $182{,}178.13$ GFLOPs/run including calibration. \\[.5em]
MSD, from scratch & 10 crops & 0,1,2; 0,2,20 & 4 models $\times$ 14 objectives; 504 runs; $31{,}976.45$ GFLOPs/run. \\[.5em]
ADNI, from scratch & 10 crops, 5 subjects & 0,1,2; 0,1,4 & 4 models $\times$ 14 objectives; 504 runs; same FLOP cap. \\[.5em]
GeoFWI3D, SAM3 & 5 scenes & 0,1,3; 0,1,2 & 14 objectives; 126 runs; 25 updates/run. \\[.5em]
MSD, MedSAM3 & 5 crops & 0,1,3; 17,83,191 & 14 objectives; 126 runs; 200 updates/run. \\[.5em]
ADNI, MedSAM3 & 2 crops, 2 subjects & 0,1,3; 0,1,2 & 14 objectives; 126 runs; 200 updates/run. \\
\hline
\end{tabular}
\caption{Completed study design: 1,890 objective runs in 210 conditions,
comprising 630 geological and 1,260 medical runs.
Each MedSAM3 campaign covers fourteen objectives, counting the baseline
objective once. Repeated matched controls (baseline runs) support paired gains but do
not enlarge the objective inventory.
Geological recipes and strengths were tuned during development;
MedSAM3 fold selection is exploratory.
All scores use matched final-budget baselines.}
\label{tab:experimental-design}
\end{table}

The medical conventional seed sets were chosen after inspecting
preliminary outcomes, so these comparisons are exploratory.
The numeric seed values identify pseudorandom streams; equal numbers
across different models or datasets would not make their initial states
identical. Each study retains its recorded seed inventory for reproducibility.
Within each dataset, all four architectures and all fourteen objectives
share the same seed set,
preserving matched comparisons between each constrained objective and
its baseline. Each objective includes all nine combinations of three
folds and three seeds, with equal weight for the three folds.

An annotated \emph{volume} is the stored example. In MSD, this is a
hippocampal crop, so a ten-volume or five-volume support set is an
annotation budget in crops. The available identifiers do not establish
how those crops map to participants. ADNI retains subject identifiers
and allows subject grouping to be checked explicitly.
Table~\ref{tab:dataset-task-overview} summarizes the images and labels.
The examples in the following sections pair an image slice with its
reference mask and a transparent overlay. They use the native grids,
before model-specific preparation, and show the slice with the largest
labelled area in the chosen viewing direction. Only MSD includes an
internal anterior/posterior division. The experiments evaluate
segmentation, with no downstream diagnostic task.

\begin{table}[htbp]
\centering
\small
\setstretch{1.05}
\renewcommand{\arraystretch}{1.2}
\begin{tabular}{P{2.65cm}P{3.15cm}P{3.65cm}P{4.1cm}}
\hline
Dataset & Model input & Reference classes & Role in the investigation \\
\hline
GeoFWI3D & Synthetic 3D seismic reflectivity &
$0$: background; $1$: salt &
Nonmedical four-model comparison and SAM3 adaptation. \\[0.4em]
MSD hippocampus & Single-channel MRI crop &
$0$: background; $1$: anterior; $2$: posterior &
Main medical comparison, including the FLOP-capped study and MedSAM3 adaptation. \\[0.4em]
ADNI hippocampus & Single-channel MRI; separate left/right crops &
$0$: background; $1$: hippocampus &
Second medical dataset for whole-hippocampus segmentation. \\
\hline
\end{tabular}
\caption{Dataset tasks and class encodings in the local preparations.
GeoFWI3D uses reflectivity as input; velocity and relative geological time
are used only to construct its reference salt mask. ADNI laterality
identifies separate crops and does not create separate foreground classes.
Experimental subsets and partitioning are specified in this chapter.}
\label{tab:dataset-task-overview}
\end{table}


\FloatBarrier
\subsection{GeoFWI3D: Salt Segmentation}
\label{sec:geo-methods}


The geological collection is drawn from archive 0 of GeoFWI3D
\cite{swaminadhan2026geofwi3d}, containing 1,000 synthetic scenes on
$96\times96\times96$ grids. The experiment retains 479 scenes with
salt: 249 without faults and 230 with faults. This selection defines
a delineation task conditional on salt being present. Entirely
salt-free scenes are outside its evaluation population.

Each example pairs a reflectivity volume with a binary annotation that
identifies salt. Figure~\ref{fig:geo-data-examples} shows one cross-section:
the image contains the measurements available to the model, and the blue
region marks the salt body it is trained to identify.

\begin{figure}[htbp]
\centering
% Native data; generator: figures/source/dataset-methods-examples.py.
% Provenance: reports/thesis_dataset_methods_examples_20261103/provenance.json.
\includegraphics[width=0.98\linewidth]{figures/compiled/geo-data-examples.pdf}
\caption{One GeoFWI3D example: reflectivity image, reference salt mask and
their overlay on the same slice. Blue denotes salt; white in the mask
denotes background. The slice is taken from a $96\times96\times96$
volume, with geological depth increasing downward. Reflectivity is
displayed over $[-1,1]$. The image illustrates the input and training
annotation, not a model prediction.}
\label{fig:geo-data-examples}
\end{figure}

The salt annotation uses two fields supplied by the simulator: velocity
and relative geological time (RGT), a coordinate describing the ordering
of geological layers. A voxel is labelled as salt when its RGT is zero
and its velocity matches the scene's recorded salt value. Using RGT alone
would also include null background artifacts. The preparation checks that
every voxel with the recorded salt velocity lies in the zero-RGT region
and retains all such components, including disconnected ones. This avoids
treating the background artifacts as salt boundaries.
% Verified label rule: scripts/geofwi3d_data.py, salt_mask.
% Velocity equality uses np.float32(metadata.salt_vel), as in the source volume.

Each reflectivity image is divided by the 99.5th percentile of its
absolute intensities and clipped to $[-1,1]$. This normalization uses
only image values and preserves the native $96^3$ grid.

For training from scratch, folds 0, 1 and 2 each supply 19 labelled
training scenes and 96 validation scenes. The support set is a fixed
approximately 5\% subset of the eligible training pool. Training and
validation are disjoint within each fold, and the three validation sets
are pairwise disjoint. The support sets contain 54 distinct scenes across
their 57 entries. All models and objectives use these same partitions
and the seeds in Table~\ref{tab:experimental-design}. SAM3 uses the smaller
support sets specified in Section~\ref{sec:sam3-methods}.

\subsection{MSD: Medical Segmentation Decathlon}
\label{sec:msd-data-methods}


The downloaded MSD hippocampus dataset used here contains 260 labelled
MRI crops with background, anterior and posterior hippocampus classes
\cite{simpson2019dataset}. All 260 image--label pairs are retained in the
local preparation and form the pool for the experiments and annotation
audit.
% Verified against Task04_Hippocampus.tar: descriptor and inventory both 260.
% Archive matches MONAI's published MD5; all local image/label bytes match.
% Audit: reports/thesis_chapter4_revision_20261103/msd_count_verification.json.
Figure~\ref{fig:msd-data-examples} shows one crop and its annotation.
Together, the blue and orange regions form the hippocampus; their shared
interface is the internal division that the model must recover in addition
to the outer boundary.

\begin{figure}[htbp]
\centering
% Native data; generator: figures/source/dataset-methods-examples.py.
% Provenance: reports/thesis_dataset_methods_examples_20261103/provenance.json.
\includegraphics[width=0.98\linewidth]{figures/compiled/msd-data-examples.pdf}
\caption{One MSD example: an MRI slice, its reference labels and their
overlay. Blue marks the anterior hippocampus and orange the posterior
hippocampus; white in the mask denotes background. All three panels show
the same slice of the native $35\times51\times35$ crop. MRI contrast is
adjusted for display using the volume's first and ninety-ninth intensity
percentiles. The coloured regions are reference annotations, not model
predictions.}
\label{fig:msd-data-examples}
\end{figure}

For the capped experiments, each image is z-score normalized: the mean of
its nonzero intensities is subtracted and the result divided by their
standard deviation. The image and its labels are then padded (extended with
zeros) and centrally cropped to $64^3$. This preserves voxel scale rather
than resizing the anatomy. Validation uses this prepared grid. Each of folds
0, 1 and 2 supplies ten support crops and 52 validation crops, shared across
models and objectives. Seeds are 0, 2 and 20. Validation identifiers are
disjoint across the three folds, but the available crop identifiers do not
establish participant-level independence. MedSAM3 uses its separate
native-grid preparation and manifests, files listing the cases assigned to
each partition, in Section~\ref{sec:medsam-methods}.

\subsubsection{Testing Proposed Anatomical Rules}
\label{sec:annotation-compatibility}

The MSD data descriptor \cite{simpson2019dataset} cites the manual
annotation protocols of Pruessner \emph{et al.}\
\cite{pruessner2000volumetry} and Woolard and Heckers
\cite{woolard2012asymmetry}. It identifies the last hippocampal-head slice
as the coronal slice containing the uncal apex, the landmark used to mark
the end of the head; body and tail form the posterior label. Coronal planes
divide the anatomy into front and back portions; sagittal planes divide left
and right. Turning this description into an automatic constraint requires a
detectable landmark, a known orientation and a relation that remains
meaningful on the stored grid. The annotation audit tests these operational
assumptions before imposing them during learning.

Two analyses address the proposed A/P transition rules. A planar analysis tests
whether one coronal separation reproduces the stored labels across the
complete 260-crop collection. Coronal-profile and image-cue analyses examine 52 original
validation cases for a reliable way to locate the transition.

The image-cue investigation includes automated and visual searches but
has no independently validated expert localization of the uncal apex.
The audit can therefore establish whether a tested numerical proxy fits
the available labels. It cannot determine whether the original annotators
misapplied the anatomical protocol. Appendix~\ref{sec:annotation-results}
reports the findings and their consequences for constraint design.


\subsection{ADNI: Alzheimer's Disease Neuroimaging Initiative}
\label{sec:adni-methods}

The ADNI dataset used here combines the final released hippocampal labels
and oriented images from the European Alzheimer's Disease Consortium
(EADC)--ADNI harmonized-protocol resource \cite{harpReleasedLabels}. Image
and label pairs are identified through subject and image metadata. Because
the release notes a possible alignment issue in converted NIfTI images, the
preparation uses the original volumetric images, stored in Medical
Imaging NetCDF (MINC) format, when exporting paired crops.

The preparation contains 135 subject identifiers and 135 image
identifiers, each contributing left and right hippocampal crops.
The resulting 270 volumes have size $32\times48\times40$ and binary
labels: zero for background and one for hippocampus. Laterality
identifies the crop; it does not define a further segmentation class.
Figure~\ref{fig:adni-data-examples} shows paired left and right crops,
their binary reference masks and annotation overlays. Both sides use
the same foreground class, without the anterior/posterior subdivision
of MSD.

% Native data; generator: figures/source/dataset-methods-examples.py.
% Provenance: reports/thesis_dataset_methods_examples_20261103/provenance.json.
\begin{figure}[htbp]
\centering
\includegraphics[width=0.98\linewidth]{figures/compiled/adni-data-examples.pdf}
\caption{Paired ADNI examples from one subject: the left and right
hippocampal crops, each shown as an MRI slice, a reference mask and an
annotation overlay. Blue denotes hippocampus on both sides; white in the
masks denotes background. Each row is a slice through a separate
$32\times48\times40$ volume. MRI contrast is adjusted separately for each
crop using its first and ninety-ninth intensity percentiles. The panels
follow the stored grids; no anatomical plane or physical scale is inferred
from the exported coordinate metadata. These are reference annotations,
not model predictions.}
\label{fig:adni-data-examples}
\end{figure}

Crops are centred on the reference-mask bounding box, with identical
indices applied to the image. Their evaluation therefore assumes prior
hippocampal localization. For the capped experiments, images are
z-score normalized over nonzero intensities and padded to $64^3$ without
resizing. The exported crops do not retain the original spacing and
orientation metadata, so their spatial constraints use voxel coordinates.

Folds 0, 1 and 2 keep both sides of each subject together. Each has an
eligible training pool of 108 subjects (216 crops) and 27 validation
subjects (54 crops), with no within-fold subject overlap. Five training
subjects contribute both sides, giving ten support crops per fold shared
across models and seeds. MedSAM3 uses the separate two-crop support
protocol in Section~\ref{sec:medsam-adni-methods}.

\subsection{Four-Model Comparisons under a FLOP Ceiling}
\label{sec:main-msd-methods}

\subsubsection{Models and Training Configurations}

The four models represent different approaches to volumetric
segmentation: a compact convolutional U-Net, the configured convolutional
system nnU-Net, UNETR with attention across image patches, and SwinUNETR
with attention within shifted local windows. Chapter~\ref{sec:background}
introduces these architectures. Testing all four examines whether the
benefit of spatial supervision depends on the model receiving it.

All start from random initialization and process one volume per update.
Within each dataset, they use the same support examples and prediction
grid: $64^3$ for MSD and ADNI, and $96^3$ for GeoFWI3D. MSD requires three
output classes; ADNI and GeoFWI3D require two. Model sizes, optimizer
settings and learning-rate schedules are recorded in
Appendix~\ref{sec:conventional-implementation}.

Table~\ref{tab:capped-baseline-losses} summarizes the medical baseline
objectives. SwinUNETR, UNETR and the compact U-Net use soft Dice over all
classes, implemented with the Medical Open Network for Artificial
Intelligence (MONAI) library. nnU-Net combines foreground Dice with
cross-entropy over all classes. On GeoFWI3D, these choices are retained
except for UNETR, which uses foreground Dice plus cross-entropy and
squared prediction terms in the Dice denominator. Its dropout is 0.1
and its initial learning rate is $6\times10^{-4}$. These differences
are held fixed within each baseline/constrained comparison.

\begin{table}[htbp]
\centering
\small
\begin{tabular}{P{4.3cm}P{3.2cm}P{6.0cm}}
\hline
Model & Baseline objective & Training Dice classes and supervision \\
\hline
SwinUNETR, UNETR and 3D U-Net
& Soft Dice
& All dataset classes, including background; final output. \\[0.5em]
nnU-Net
& Soft Dice plus cross-entropy
& Foreground classes in Dice; all classes in cross-entropy;
  multiple output resolutions. \\
\hline
\end{tabular}
\caption{Baseline losses in the capped MSD and ADNI comparisons. The class set
listed for Dice concerns training, not the selection of evaluation metrics.}
\label{tab:capped-baseline-losses}
\end{table}

nnU-Net applies its baseline losses at several output resolutions,
resampling the reference labels to each grid. All auxiliary terms act
only on the primary, full-resolution prediction. Each model also keeps
the same learning-rate trajectory across its objectives. If an extra
loss leaves room for fewer updates under the cap, that trajectory is
not shortened to compensate.

The final geological UNETR configuration includes two changes developed
after its initial recipe struggled within the budget. First, training
uses right-angle rotations and reflections in the horizontal plane,
applied jointly to images and labels. These transformations preserve
the depth axis and require no interpolation. Second, the outer two
voxel layers are excluded from band and edge supervision. A salt body
can continue beyond the observed cube, so its intersection with a cube
face need not be an observed salt/background boundary. The full input,
prediction grid and baseline loss are retained. Both choices apply to
the UNETR baseline and its constrained variants; the other three
geological models retain their original crop-face convention.
Appendix~\ref{sec:geo-boundaries} specifies that convention, and
Appendix~\ref{sec:geo-development} records the development comparison.
Thus, cross-model differences concern complete training configurations,
including these choices, rather than architecture alone.

\subsubsection{Spatial Objectives}

Fourteen objectives cover the combinations in
Table~\ref{tab:constraint-combinations}; $\varnothing$ denotes no
selected term among $B,E,S$. These abbreviations refer to boundary
bands, edge consistency and signed distance; $D$ adds depth supervision.
Preliminary MSD screening did not establish a consistent incremental
benefit from edge alone over the supervised baseline, or from edge plus
signed distance over signed distance alone. These two
configurations were therefore omitted from the subsequent
fourteen-objective replications across datasets. The study therefore
compares the retained objectives; it does not test every possible
combination or interaction.
% EXCLUSION EVIDENCE: reports/compute_matched_constraints_20pct_20261004/
% summary.json (U-Net, MSD fold 0 / seed 0): E 41.33 vs baseline 42.14;
% E+S 57.25 vs S 61.15, although E+S exceeds baseline. The longer pilot in
% reports/compute_matched_constraints_682338/README.md records only a
% +0.055 pp validation-selected edge gain and a slightly lower final score.
% These archives do not establish ineffectiveness across all datasets/seeds.
% Pilot records are preserved; omitted configurations were not part of the
% later fourteen-objective replication inventory.

\begin{table}[htbp]
\centering
\small
\setstretch{1.05}
\renewcommand{\arraystretch}{1.2}
\begin{tabular}{@{}P{4.6cm}P{4.6cm}P{4.6cm}@{}}
\hline
Subset of $B,E,S$ & Without depth & With depth \\
\hline
$\varnothing$ & Baseline & Baseline $+D$ \\
$B$ & Baseline $+B$ & Baseline $+B+D$ \\
$E$ & Not replicated & Baseline $+E+D$ \\
$S$ & Baseline $+S$ & Baseline $+S+D$ \\
$B,E$ & Baseline $+B+E$ & Baseline $+B+E+D$ \\
$B,S$ & Baseline $+B+S$ & Baseline $+B+S+D$ \\
$E,S$ & Not replicated & Baseline $+E+S+D$ \\
$B,E,S$ & Baseline $+B+E+S$ & Baseline $+B+E+S+D$ \\
\hline
\end{tabular}
\caption{The fourteen objective combinations in each conventional dataset
comparison and in the SAM3 and MedSAM3 adaptation campaigns.
Each included entry retains the architecture's supervised
baseline. Coefficients are calibrated for the corresponding combination;
adding a family need not leave every existing coefficient unchanged.}
\label{tab:constraint-combinations}
\end{table}

Bands and edges supervise the combined anterior/posterior foreground
on MSD, the hippocampus on ADNI and salt on GeoFWI3D. Signed-distance
fields (SDFs) are constructed separately for each foreground class;
depth supervision retains the original reference classes. All
conventional models use two-step bands with six face-sharing neighbours.
The class groups determine which interfaces receive supervision, while
the geological UNETR exclusion changes where bands and edges apply.

Distance fields are clipped at eight voxels and normalized by eight.
Depth supervision balances class errors across layers two units thick
within the foreground and an eight-unit surrounding region. On MSD,
these units are millimetres; its verified isotropic spacing of
$1\,\mathrm{mm}$ makes them numerically equal to voxel units. ADNI and
GeoFWI3D use voxel units and orient the reference plane using the
image grid rather than an anatomical direction. The plane passes
through the foreground centre and contains its main axis, following
Chapter~\ref{sec:spatial-constraints}. Matching these numerical widths
does not establish equal physical distances across datasets.

For GeoFWI3D, source axis 0 supplies the plane's reference direction;
another grid axis is used if its projection is too small. UNETR's
coordinate mapping is restored after augmentation before computing
the depth weights. This annotation-derived coordinate is not a
measurement of geological age or a guaranteed vertical ordering of
layers. Numerical distance and boundary conventions are given in
Appendix~\ref{sec:sdf-construction-details}.
% Geo depth settings retained from the verified protocol: projection norm
% below 0.1 selects the least-parallel source axis, with axis-index tie break;
% sparse-stratum threshold 64 voxels. No inference of physical spacing.

\subsubsection{Calibration and Activation}
\label{sec:calibration-methods}

An auxiliary loss can dominate training simply because its numerical
scale differs from the baseline. Calibration chooses its coefficient
by comparing their gradients on the support examples, without changing
the network during this measurement. The resulting coefficients remain
fixed during training. All ten medical support crops, or all 19
geological support scenes, contribute to calibration.

The conventional models first train with the baseline alone for nine
updates in nnU-Net or 50 in the other architectures. Selected terms are
then calibrated in order: bands, edges, signed distance and depth.
Each is compared with the baseline and any preceding terms, so its
coefficient can change between combinations. Targets for the typical
and unusually large gradient contributions control the strength of
supervision. Their values differ by model and dataset; the geological
nnU-Net and UNETR settings followed exploratory comparisons.
Appendix~\ref{sec:detailed-calibration} gives the equations and settings.

After calibration, auxiliary contributions increase gradually to full
strength. Signed distance either enters with the other terms or starts
later, following Table~\ref{tab:sdf-activation-policies}. Combinations
containing both signed distance and depth use the later start. These
policies were developed from earlier MSD validation experiments and
retained for the reported campaigns. The comparisons therefore assess
complete calibrated objectives, including their activation schedules.

\subsubsection{Resource Accounting and Evaluation}

A fixed computation allowance limits how much training each model can
receive. The medical ceiling is approximately 31,976 GFLOPs, the
measured cost of 90 nnU-Net updates. With ten support crops, this
reference represents nine passes through the training subset.
Exploratory learning curves informed this short allowance. Models
with cheaper updates can take more steps within it.

The geological study uses a larger allowance for its larger input
volumes: approximately 182,178 GFLOPs, defined from the original cost
of 152 nnU-Net updates, or eight passes through 19 support scenes.
The final training configurations retain that ceiling. Medical
calibration is recorded separately, whereas geological calibration
consumes part of the allowance. The budgets therefore differ in both
size and accounting scope across domains.

Training costs include the forward pass, active losses, backpropagation
and optimizer update. Another step is taken only if its full cost fits
in the remaining budget. Reference-target construction, validation,
data loading, storage and recipe development are outside these
allowances. The comparison matches the specified arithmetic work,
not elapsed time or end-to-end computation, and does not imply that
different architectures are equally close to convergence.
Figure~\ref{fig:flop-budget-example} illustrates these distinctions.

\begin{figure}[htbp]
\centering
% Generator: figures/source/explanatory-concepts.py
\includegraphics[width=\linewidth]{figures/compiled/flop-budget-example.pdf}
\caption{Computation budgets in the comparisons from scratch.
(a) Each block bounded by white separators represents one complete
optimizer update. A more expensive update permits fewer steps under
the same ceiling; auxiliary computation also consumes that allowance.
Costs and update counts here are illustrative, not experimental
measurements. (b) Medical comparisons report calibration separately
from their optimization ceiling. The geological ceiling includes both
calibration and optimization. The outlined containers indicate
accounting scope, not relative measured costs or equal budgets across
datasets.}
\label{fig:flop-budget-example}
\end{figure}
\FloatBarrier

Each of the three conventional studies contributes all 504
final-budget runs. Reporting the highest-mean objective afterwards
is a descriptive comparison, separate from checkpoint selection.
Appendix~\ref{sec:detailed-resource-accounting} gives the exact ceilings
and counting conventions. Table~\ref{tab:cap-training-exposure} shows
the resulting numbers of updates and passes through the support data.

\subsection{Adapting Pretrained Models}

The adaptation studies ask whether spatial supervision also helps
when the model already has pretrained image representations. SAM3 and
MedSAM3 process two-dimensional image planes, although the annotations
and final evaluation concern three-dimensional volumes. During
training, neighbouring planes are sampled together so that the edge
term can compare predictions across slices as well as within them.
At evaluation, slice predictions are reassembled into a volume.

All three studies freeze the SAM3 base weights and update only the
low-rank adaptation (LoRA) factors introduced in Chapter~\ref{sec:background},
giving 18,497,664 trainable parameters. Geological adaptation starts
with fresh factors that initially leave the base model unchanged;
medical adaptation starts with the released MedSAM3 adapter
\cite{carion2025sam3,liu2025medsam3}. Every matched baseline and
constrained run restores the same base weights and initial factors,
and uses the same support images and sampling schedule.

The pretrained baseline supervises object classification, bounding
boxes, target presence and segmentation masks. Its mask losses combine
soft Dice with focal loss, a form of cross-entropy that gives less
weight to already confident correct predictions. Spatial penalties
are added to this objective and calibrated on the support annotations.
Each update uses two adjacent-slice pairs. Constraint regions and
normalization counts are defined on the full annotated volume, so
sampling preserves the intended weighting across it; the corresponding
reductions are given in Appendix~\ref{sec:sampled-pretrained-supervision}.

Inference uses image planes and fixed text queries. The model combines
its confidence in each mask location, its match to the query and the
target's presence to form foreground responses. These responses are
converted into mutually exclusive class labels. Reference masks are
used to score the predictions, never to supply localization prompts.

\paragraph{SAM3 on GeoFWI3D.}\label{sec:sam3-methods}
Each fold uses five labelled support volumes, taken from the beginning
of its established support list, and 96 held-out volumes. The native
$96^3$ reflectivity volumes are mapped from $[-1,1]$ to $[0,1]$ for
image processing. The query \emph{salt body} produces one foreground
response; its complement represents background.

Every objective receives 25 updates, using the first 25 steps of the
retained 200-update learning-rate schedule. The schedule is not compressed
into a shorter decay. Bands are calibrated at initialization; other selected
terms are calibrated after five prefix updates, during which their weights
are zero. Bands and edges use the salt mask, with two-step, zero-padded
morphology retaining crop-face support. SDF uses one foreground field on the
unit voxel grid. Depth weights use source-array geometry, with no anatomical
orientation assumption. These choices distinguish this adaptation protocol
from the conventional geological recipes.

\paragraph{MedSAM3 on MSD.}\label{sec:medsam-methods}
Five support crops are used in each fold, with 46, 52 and 44 held-out
crops respectively. The 15 distinct support crops and 142 distinct
evaluation crops do not overlap. Fold selection followed earlier
exploratory outcomes, and participant identities are unavailable; the
crop separation therefore does not establish participant independence.
Images and masks are reordered and, where needed, flipped together
into the closest right--anterior--superior orientation recorded by their
NIfTI coordinate transforms. Intensities are clipped and scaled between the first and
99th percentiles, then converted to eight-bit planes.

The fixed queries are \emph{anterior hippocampus} and \emph{posterior
hippocampus}. Each response probability $q$ is converted to log-odds,
$\log(q/(1-q))$. These scores, together with a zero background logit, are
normalized by softmax before argmax assigns one class per voxel. Bands and
edges are applied separately to background, anterior and posterior, so they
include the internal interface. SDF uses the two foreground classes and
depth uses physical geometry. All variants receive 200 updates, with their
selected auxiliary terms ramping from the first update.

The fourteen objectives were completed across two execution batches.
The original bands and bands-plus-edge results retain their matched
baseline and calibration; the other eleven constrained variants use
the controls and calibration recorded with their batch. Those controls
are not pooled, and repeated baseline runs do not create additional
objective configurations.
Figure~\ref{fig:medsam3-workflow} shows the adaptation and prediction
path.

\begin{figure}[htbp]
\centering
\setstretch{1.05}
\begin{tikzpicture}[x=1cm,y=-1cm,font=\footnotesize,
  every node/.style={text=thesisInk},>=stealth]
\node[align=center,text width=15cm,inner sep=0pt,text=thesisGray,font=\fontsize{9}{11}\selectfont\bfseries] at (7.750,0.200) {TASK-SPECIFIC ADAPTATION ON MSD};
\draw[fill=thesisPaleGray,draw=thesisGray,line width=.65pt,rounded corners=2pt] (0.150,0.625) rectangle (7.450,1.775);
\node[align=center,text width=7.06cm,inner sep=0pt,text=thesisInk,font=\fontsize{10}{12}\selectfont] at (3.800,1.200) {SAM3 base weights\\Frozen throughout adaptation};
\draw[fill=thesisPaleOrange,draw=thesisOrange,line width=.65pt,rounded corners=2pt] (8.050,0.625) rectangle (15.350,1.775);
\node[align=center,text width=7.06cm,inner sep=0pt,text=thesisInk,font=\fontsize{10}{12}\selectfont] at (11.700,1.200) {Released medical LoRA adapter\\Starting factors restored for each arm};
\draw[draw=thesisGray,line width=0.8pt,->] (3.800,1.775) -- (3.800,2.150) -- (7.750,2.150) -- (7.750,2.500);
\draw[draw=thesisGray,line width=0.8pt] (11.700,1.775) -- (11.700,2.150) -- (7.750,2.150);
\draw[fill=thesisPaleBlue,draw=thesisBlue,line width=.65pt,rounded corners=2pt] (0.150,2.500) rectangle (15.350,3.600);
\node[align=center,text width=14.959999999999999cm,inner sep=0pt,text=thesisInk,font=\fontsize{10}{12}\selectfont] at (7.750,3.050) {MedSAM3 initialization\\Pretrained base + medical adapter};
\draw[draw=thesisGray,line width=0.8pt,->] (11.700,3.600) -- (11.700,4.150);
\draw[fill=thesisPaleBlue,draw=thesisBlue,line width=.65pt,rounded corners=2pt] (0.150,4.150) rectangle (7.150,6.250);
\node[align=center,text width=6.76cm,inner sep=0pt,text=thesisInk,font=\fontsize{10}{12}\selectfont] at (3.650,5.200) {5 labelled support volumes\\2 adjacent-slice pairs per update\\Fixed anterior / posterior queries};
\draw[fill=thesisPaleOrange,draw=thesisOrange,line width=.65pt,rounded corners=2pt] (8.050,4.150) rectangle (15.350,6.250);
\node[align=center,text width=7.06cm,inner sep=0pt,text=thesisInk,font=\fontsize{10}{12}\selectfont] at (11.700,5.200) {Baseline / baseline + constraints\\200 updates; only LoRA factors (18.50 M)\\14 objectives; two execution batches};
\draw[draw=thesisGray,line width=0.8pt,->] (7.150,5.200) -- (8.050,5.200);
\draw[draw=thesisGray,line width=0.3pt] (0.100,6.700) -- (15.400,6.700);
\node[align=center,text width=9.5cm,inner sep=0pt,text=thesisGray,font=\fontsize{9}{11}\selectfont\bfseries] at (5.000,7.150) {INFERENCE ON HELD-OUT VOLUMES};
\draw[draw=thesisGray,line width=0.8pt,->,dashed] (11.700,6.250) -- (11.700,7.800);
\draw[fill=thesisPaleBlue,draw=thesisBlue,line width=.65pt,rounded corners=2pt] (0.150,7.800) rectangle (7.150,9.100);
\node[align=center,text width=6.76cm,inner sep=0pt,text=thesisInk,font=\fontsize{10}{12}\selectfont] at (3.650,8.450) {MRI slices\\Fixed anterior / posterior queries};
\draw[fill=thesisPaleBlue,draw=thesisBlue,line width=.65pt,rounded corners=2pt] (8.050,7.800) rectangle (15.350,9.100);
\node[align=center,text width=7.06cm,inner sep=0pt,text=thesisInk,font=\fontsize{10}{12}\selectfont] at (11.700,8.450) {Adapted MedSAM3\\Two foreground response maps};
\draw[draw=thesisGray,line width=0.8pt,->] (7.150,8.450) -- (8.050,8.450);
\draw[draw=thesisGray,line width=0.8pt,->] (11.700,9.100) -- (11.700,9.800);
\draw[fill=thesisPaleBlue,draw=thesisBlue,line width=.65pt,rounded corners=2pt] (8.050,9.800) rectangle (15.350,11.400);
\node[align=center,text width=7.06cm,inner sep=0pt,text=thesisInk,font=\fontsize{10}{12}\selectfont] at (11.700,10.600) {Three-class semantic projection\\Background + anterior + posterior\\Voxelwise argmax};
\draw[fill=thesisPaleGray,draw=thesisGray,line width=.65pt,rounded corners=2pt] (0.150,9.800) rectangle (7.150,11.400);
\node[align=center,text width=6.76cm,inner sep=0pt,text=thesisInk,font=\fontsize{10}{12}\selectfont] at (3.650,10.600) {Reassembled 3D label map\\0: background; 1: anterior; 2: posterior};
\draw[draw=thesisGray,line width=0.8pt,->] (8.050,10.600) -- (7.150,10.600);
\node[align=center,text width=15.3cm,inner sep=0pt,text=thesisGray,font=\fontsize{9}{11}\selectfont] at (7.750,11.850) {Evaluation: Dice of final predictions; paired gains compare each variant with its baseline.};
\end{tikzpicture}
\caption{MedSAM3 adaptation and inference on MSD. Training changes only
the medical LoRA factors using five support crops per fold. The two
fixed queries produce foreground responses, which are combined with
background to obtain mutually exclusive labels. Held-out masks are used
for evaluation and supply no localization prompts.}
\label{fig:medsam3-workflow}
\end{figure}

\paragraph{MedSAM3 on ADNI.}\label{sec:medsam-adni-methods}
Each fold uses two labelled crops from two subjects and 54 held-out
crops from 27 other subjects. Folds 0 and 3 share their two support
crops; fold 1 shares one and changes the other. Only three distinct
support crops occur across the study. These folds vary the evaluation
partition, but do not provide three independent draws of a training set.

The medical initialization, 200-update schedule and intensity scaling
follow the MSD protocol. ADNI retains its stored array orientation and
uses the single query \emph{hippocampus}. The binary projection combines
the foreground response with its complement. Bands and edges supervise
the hippocampus mask, without a second background-mask average.
SDF uses unit array voxels, and depth uses source-axis geometry without
an inferred anatomical orientation.

These update budgets describe additional task adaptation and exclude
pretraining. They do not provide a FLOP match with the conventional
studies. Exact optimizer settings, calibration rules, fold and seed
identifiers, and execution environments are retained in
Appendix~\ref{sec:adaptation-implementation}. Runtime comparisons keep
their recorded scope; a training-loop timer is not an end-to-end
adaptation cost.

\subsection{Evaluation Measures}
\label{sec:evaluation-methods}

Every reported Dice score measures overlap between a discrete
prediction and its reference annotation. This is hard Dice, distinct
from the differentiable soft Dice loss used during training.
Validation scores use held-out examples; training scores use the
support examples supplied during fitting.

Each voxel receives the class with the highest predicted probability
(argmax). For volume $i$ and class $c$, let $Y_{i,c}$ be the set of
reference voxels and $\widehat Y_{i,c}$ the predicted set. Hard Dice is
\begin{equation}
 D_{i,c}=\frac{2|Y_{i,c}\cap\widehat Y_{i,c}|}
 {|Y_{i,c}|+|\widehat Y_{i,c}|}.
\label{eq:reported-hard-dice}
\end{equation}
A present reference class with no predicted voxels scores zero.
The MSD volume score is the mean anterior/posterior Dice,
$D_i=(D_{i,A}+D_{i,P})/2$, and the ADNI volume score is hippocampus Dice. Geological evaluation
uses salt Dice.
Background is excluded from these primary metrics.
Section~\ref{sec:background-dice-results} also reports background and
all-class Dice on the same predictions.

To summarize a study, we first average Dice equally across validation
volumes in a run, then across the three seeds within each fold, and
finally across the three folds. The reported sample standard deviation
describes variation between those three fold averages. Training Dice
follows the same hierarchy using support volumes.

Let $\overline D_{m,j,f,s}$ denote the mean volume Dice in a run with
model $m$, objective $j$, fold $f$ and seed $s$. Subscript
$\mathrm{base}$ identifies its matched baseline. Their paired change is
\begin{equation}
 \Delta_{m,j,f,s}=100\left(
 \overline D_{m,j,f,s}-\overline D_{m,\mathrm{base},f,s}\right).
\label{eq:paired-reported-gain}
\end{equation}
Dice is reported as a percentage and gains as percentage-point
differences. Paired gains follow the same averaging hierarchy, with
their standard deviation calculated across the three fold-mean
differences. Positive-change counts refer to the nine individual
matched runs. These standard deviations describe observed variation,
not confidence intervals over independent datasets: seeds share data,
training pools overlap and configuration selection is exploratory.
Any dispersion across individual volumes is identified separately.

Error counts complement Dice by distinguishing missed foreground
(false negatives, FN) from extra foreground (false positives, FP).
MSD also permits anterior-to-posterior and posterior-to-anterior swaps.
These four disjoint error categories account for every incorrectly
labelled voxel. They can be read from a confusion matrix, which records
reference classes against predicted classes; the pretrained-model
matrices use reference classes as rows. ADNI and GeoFWI3D are binary
tasks and have no anterior/posterior swaps.

All counts cover the full prediction grid, including errors outside
the boundary bands. Conventional studies sum errors across validation
volumes before averaging runs. MedSAM3 and SAM3 instead average errors
per validation volume, then across seeds and folds. Their tables state
these units because totals over a validation set cannot be compared
directly with averages per volume.

\clearpage

